\section{Conclusion}
\label{sec:conclusion}

In this paper, we have deconstructed the ``Re-Quantization Silence,'' exposing a widespread methodological blindspot in the model merging literature. By assuming a high-precision abstraction, previous studies have ignored the downstream compression steps required for hardware deployment. Our rigorous, multi-axial auditing framework (RQA) demonstrates that post-hoc low-bit quantization (e.g., 4-bit) of merged models catastrophically crushes task-specific updates under per-tensor grids, but is nearly lossless under standard per-channel configurations.

By introducing and evaluating Scale-Adaptive Weight Shifting (SAWS) and Quantization-Aware Adapter Coefficient Search (QA-ACS), we have laid bare the empirical and mathematical realities of both data-free scaling and optimization-based mitigations. We exposed a fundamental scale preservation dilemma in SAWS, showing that its success is driven by selective task-vector boosting, and we documented the severe fragility of unsupervised test-time optimization (QA-ACS), which collapses under high quantization noise.

\subsection{Limitations and Future Work}
We must explicitly acknowledge two major limitations of this study. First, our evaluation is restricted to a **small-scale Vision Transformer backbone (ViT-Tiny, 5.7M parameters)**. While this enables rapid multi-axial profiling, quantization and merging dynamics on toy networks may not fully generalize to the multi-billion parameter architectures (such as LLaMA or Mistral) where QLoRA and model merging are deployed in practice. Crucially, as a preliminary empirical indicator at larger scales, we evaluated the double-quantization format-shift noise on the larger \texttt{vit\_base} (86M parameters) backbone in Table~\ref{tab:double_quantization_error} (Section 3.2.1), finding that the Relative Frobenius reconstruction error remains substantial and indeed intensifies up to $32.27\%$ under INT4 configurations. However, evaluating our complete multi-task auditing and mitigation framework on multi-billion parameter models represents an urgent direction for future work\footnote{As of this writing, a pilot scaling evaluation of SAWS on Pythia-1B and LLaMA-1B under 4-bit block-wise quantization configurations is actively running on our compute cluster to empirically validate these large-scale scaling hypotheses.}.

Second, the continuous, full-precision baseline (Naive FP16 Merge) already exhibits severe performance collapse, achieving only $66.65\%$ mean accuracy compared to the unmerged expert ceiling of $93.85\%$ (with MNIST dropping to $45.40\%$). This shows that \textbf{task interference} in weight space is the primary bottleneck of this setup. Studying the quantization of a merged model that is already severely degraded in full precision introduces confounding noise into the empirical results. To decouple this task-interference bottleneck from post-training quantization, we propose that future research adopt a **Zero-Interference RQA Protocol**. Under this protocol, multiple task-specific adapters are trained on identical data distributions with non-overlapping, orthogonal head routing, or on highly domain-aligned tasks—such as multi-lingual machine translation (e.g., merging French-to-English and Spanish-to-English adapters trained on the Europarl corpus) or localized text classification (e.g., merging Amazon review sentiment classification adapters across different product categories). Because the underlying expert tasks share highly compatible representational structures, the full-precision merging process is mathematically lossless, achieving close to $100\%$ task arithmetic retention. Applying our auditing framework to this zero-interference regime would isolate the downstream discretization noise from representation conflicts, allowing an entirely pure and unconfounded characterization of re-quantization erasure. Future studies should leverage this proposed protocol to benchmark mitigations under pristine baseline conditions.

Ultimately, our findings call for a paradigm shift in the evaluation protocols of model merging and PEFT research. To build truly deployable deep learning systems, we must evaluate model fusion through the lens of hardware constraints, and future research must prioritize transparency and self-criticism above reporting inflated full-precision SOTA metrics.